\documentclass[UTF8]{ctexart}
\usepackage{geometry}
\usepackage{hyperref}
\usepackage{listings}
\usepackage{xcolor}

\geometry{a4paper, left=2.5cm, right=2.5cm, top=2.5cm, bottom=2.5cm}

\title{Agent开发周报 260401}
\author{刘德辰}
\date{2026年4月1日}

\begin{document}

\maketitle

\section{车辆报告链路继续完善：数据源切换、管道修复与指标补齐}

\textbf{背景}：上周车辆报告虽然已经接入统一报告查询方向，但车辆这一类报告本身还有一组独立问题没有解决完：一是数据源还没有完全切到 MCP 侧，二是结构化查找、脚本文档和实际输出之间还有不一致，三是车辆报告的指标下钻能力还不够完整。

\textbf{方案说明}：本周围绕车辆报告这条链路继续单独推进，先把数据源切到 MCP，再修正中间管道以及脚本文档和实际输出之间的不一致，最后补齐报告里的指标展开能力，让车辆报告本身先稳定下来，而不是和车辆专家实验混在一起推进。

\textbf{功能说明}：
\begin{itemize}
\item 将车辆报告查询切到 MCP 数据源，补上车辆资料读取能力
\item 修复车辆报告 MCP 管道中的结构化查找、路由衔接和执行链路问题
\item 统一结构化报告数据源入口，减少不同报告类型各走一套来源判断的情况
\item 继续对齐车辆报告输出与脚本文档，减少“脚本怎么写”和“结果怎么出”不一致的问题
\item 扩展车辆报告指标下钻能力，补齐部分指标只能总览、不能进一步展开的问题
\item 收紧报告 worker 的 negative prompt，并澄清报告类工具描述，减少报告生成时被错误引导或误触发
\end{itemize}

\section{车辆专家实验进入新阶段：从专家拆分到 CoT 模式对比}

\textbf{背景}：上周新的实验重点是验证“车辆问题是否值得从 omni 中拆出来交给车辆专家”，这一轮实验已经进入真实流量。本周这条线继续推进，但重点已经不是报告链路，而是专家实验本身的路由边界、回答深度和 CoT 模式对比。

\textbf{方案说明}：本周先补齐车辆专家技能入口和响应深度控制，再把实验正式推进到 CoT 模式 AB，对比“车辆专家普通模式”和“车辆专家 CoT 模式”的实际效果。同时继续收紧 omni 在车辆场景下的提示词边界，避免实验流量被通用链路提前截走。

\textbf{功能说明}：
\begin{itemize}
\item 暴露车辆专家 router tool，让车辆专家实验链路更完整、更可控
\item 为车辆专家技能补充 response-depth 路由能力，区分简短回答与展开分析两类场景
\item 正式启动车辆专家 CoT AB，并补齐相关实验配置与说明材料
\item 配套调整前端实验 UI、统计展示和联调脚本，方便继续观察实验命中与回答差异
\item 新增车辆专家 cot mode routing，让实验组可以按模式稳定落到对应回答链路
\item 补充 Y 组专用 omni negative prompting，减少 omni 对车辆问题的抢答倾向（其实是 router 分配问题）
\end{itemize}

\section{澄清与路由中间态继续统一：让补问、规则回复和路由判断走同一套入口}

\textbf{背景}：除了报告和专家两条业务线之外，这周还有一条独立工作线是“中间态治理”。之前系统里对澄清补问、进一步补信息、规则回复和路由契约存在多套处理方式，有的直接写在 Runtime 里，有的在 Router 里兜底，有的在消息拼装时特殊处理，导致多轮对话行为不够一致。

\textbf{方案说明}：本周把这部分作为单独事项推进。先抽出 clarification state helper 和通用澄清工具，再逐步把 request-further-info、rule-reply contract 和 pending clarification 都改为由 Router 统一判断，减少 Runtime 直接特判的情况。

\textbf{功能说明}：
\begin{itemize}
\item 新增 clarification state helper，统一澄清状态的记录和恢复方式
\item 将通用 clarification tool 接入 Runtime，补齐澄清能力的运行时入口
\item 收紧 rule-reply routing contract，减少规则回复与普通咨询之间互相走错路径的情况
\item 简化 request-further-info 消息管道，减少中间态消息在多处重复拼装
\item 将 pending clarification 正式改为交由 Router 继续判断，而不是由 Runtime 直接分支处理
\item 继续调整 rule-threshold router prompt，降低过于生硬的阈值判断带来的误分流
\end{itemize}

\section{RAG 服务从“可用”推进到“可接入”：上线交接、切块优化与检索增强}

\textbf{背景}：知识库/RAG 这周也是一条完全独立的工作线。前期主要完成了基础服务搭建和入库链路，本周开始处理两类问题：一类是服务层面的上线交接、测试服恢复和文档补齐；另一类是检索质量本身，包括切块结构、召回效果和重排能力。

\textbf{方案说明}：本周分两段推进。先把 RAG 服务当前状态、测试服务器恢复和交接流程整理清楚，明确服务已进入可用状态；再围绕层级切块、query rewrite 和 coarse rerank 继续优化检索质量，为后面接入 omni 提供更稳定的上下文来源。

\textbf{功能说明}：
\begin{itemize}
\item 明确 RAG 服务已可用，并补齐测试服务器恢复、上线交接和部署说明文档
\item 优化文档层级切块逻辑，改善 heading path、块合并和层次信息保留效果
\item 补充 document parser、job worker 和知识库服务相关测试，降低后续继续迭代切块逻辑时的回归风险
\item 在知识库详情页展示合并后的 chunk 细节，方便直接检查实际切块结果
\item 新增 query rewrite、LLM 粗排和 rerank 能力，继续提升检索召回和排序质量
\item 将知识库上下文正式注入 omni，为后续问答链路使用 RAG 提供入口
\end{itemize}

\section{分析页前端重建：主体迁到 Vue，并补齐构建与本地化}

\textbf{背景}：分析页的前端改造是这周体量最大的一条工作线，而且和 Agent 主链路、RAG、车辆报告都没有直接关系。此前分析页长期堆在 React/TSX 结构里，页面重、图表多、维护成本高，继续在旧结构上加内容已经不太合适。

\textbf{方案说明}：本周直接把分析页主体按 Vue 方案重建，并把旧 TSX 实现迁到 legacy 目录保留参考。迁移完成后，再补锁文件、构建依赖和中文文案，使新页面先达到“能跑、能看、能继续迭代”的状态。

\textbf{功能说明}：
\begin{itemize}
\item 重建分析页工作台、图表注册体系和多类分析页面，主体切换到 Vue 实现
\item 新增 Vue 版分析入口、宿主页和图表封装，形成新的承载结构
\item 将旧的 TSX 分析页整体迁入 legacy，方便迁移期对照和必要时回退
\item 补齐 Vue 分析页相关依赖和 lockfile，确保新方案可以稳定构建
\item 补做分析页中文文案和本地化数据处理，减少迁移后界面表述不统一的问题
\item 修复助手页 streaming session sync 问题，避免前端会话展示出现错位
\end{itemize}

\section{Runtime 拆分与文档补齐继续推进：减少主文件堆积，保持结构可追踪}

\textbf{背景}：运行时主文件这几周一直在持续瘦身。本周虽然没有像分析页迁移那样大面积改造，但仍有两类重要工作：一是继续把 Runtime 里的研究接线和报告相关逻辑往外迁；二是同步更新文档、实验说明和周边材料，避免结构已经变化但文档还停留在旧版本。

\textbf{方案说明}：本周按“先迁高频改动区域，再同步快照文档”的顺序继续推进。把 Runtime 中与 research wiring、structured report 相关的部分继续剥离，再把最新结构、实验说明和周报材料同步写入文档。

\textbf{功能说明}：
\begin{itemize}
\item 继续瘦身 Runtime，剥离 research wiring，减少主文件堆积
\item 新增 \texttt{structured-report-runtime}，把报告相关运行时逻辑独立成模块
\item 更新 \texttt{runtime.md}，同步最新拆分边界和运行时快照
\item 补充与整理实验、MCP、运行时结构和交接相关文档
\item 提交并归档上一版周报等阶段性材料，降低后续复盘成本
\end{itemize}

\section{技术债务与限制}

\begin{itemize}
\item \textbf{车辆报告链路虽然继续完善，但还需要继续回归}：数据源已经切到 MCP，脚本文档和指标也在补齐，但车辆报告的稳定性还需要更多真实样例验证。MCP 交付文档需要包含验收标准和测试用例。
\item \textbf{报告追问与多轮对话还需要继续验证}：虽然这周把澄清中间态进一步统一到 Router，但报告生成后的追问、补问和业务咨询之间能否长期稳定衔接，还需要继续回归
\item \textbf{车辆专家实验仍在观察期}：本周已经推进到 CoT 模式对比，但“是否该拆专家”以及“CoT 是否真的有收益”都还没有最终结论
\item \textbf{澄清中间态刚统一回 Router}：pending clarification 和 request-further-info 已经开始统一，但复杂多轮对话下是否稳定，还需要继续测
\item \textbf{MCP 对接仍受外部服务质量影响}：虽然这周已经把车辆报告数据源切到 MCP，但数据没有完全接入，接口行为测试比较受局限
\item \textbf{RAG 已经进入可接入阶段，但检索质量还在优化中}：切块、rewrite 和 rerank 已经补上，真实问答中的收益和误召回情况还要继续评估
\item \textbf{分析页 Vue 迁移完成的是主体，不是全部收尾}：本周先把大框架切过去，后续还需要继续补细节、做联调和回归
\item \textbf{Runtime 拆分仍未完成}：虽然主文件继续变薄，但还有一部分上下文管理和业务接线逻辑没有完全迁出
\item \textbf{规则模板库}：建立常用规则模板，降低配置成本（无限期延后）。
\item \textbf{本地模型性能估算还没有启动}：还没有基于 token 消耗量去估算本地部署模型后的响应速度、吞吐和资源需求。这件事和当前几条主线相对独立，但后面需要专门安排时间做一次完整评估，避免本地部署阶段再临时补课
\end{itemize}

\section{下周计划}

\begin{enumerate}
\item \textbf{继续回归车辆报告链路}：重点验证 MCP 数据源切换后的车辆报告稳定性、指标下钻以及脚本文档与实际输出的一致性
\item \textbf{继续观察车辆专家 CoT 实验数据}：结合统计与真实会话判断 CoT 模式是否值得继续推进
\item \textbf{补齐澄清链路多轮回归}：重点看 pending clarification、rule-reply 和补问链路在复杂上下文下是否稳定
\item \textbf{继续优化 RAG 接入效果}：围绕切块、rewrite、rerank 和 omni 上下文注入继续调优
\item \textbf{继续处理分析页 Vue 迁移后的收尾问题}：补联调、补细节、补回归
\item \textbf{继续推进 Runtime 拆分与文档同步}：减少代码结构和文档描述之间的偏差
\item \textbf{准备 MCP 服务的后备方案}：如果当前 MCP 服务无法满足需求，需要准备后备方案。
\end{enumerate}

\end{document}
